\cotefile{84}
\scanpage{82}


Donc la cat. des \unsure{-torseurs|torseurs} $2$-trivialisés, pour des Modules sur lesquels $2$ est inv., équivaut à celle des Modules sans plus.

On peut, sans \unsure{restr.} sur $2$ (\unsure{i.e.|inv.} sur $X$) ou sur $M$, regarder les \deleted{modules} torseurs $2$-trivialisés qui ont cette structure : munis de \inserted{scindages} $x_0\in\Gamma P_1$,
tels que $x_0=\overline{x_0}$, d'où $E\simeq L_0\oplus\underline{O}$, $T\simeq(0,2)$, comme dessus. On aura.

\[ b_{x_0,T}=2x_0-T=0,\]

La théorie est \unsure{tot.!|t.t.|tt. !} triviale \deleted{pour} $2$ inv. sur $X$ ! Le produit $*$ se réduit au produit $\otimes$ des \unsure{$L_i$|$L^{\mathrm i}$} sans plus, ``les $b$ suivent'' -- ils sont nuls !

Le cas \inserted{le plus \deleted{intéressant}} intéressant est celui où $2$ n'est \deleted{\unsure{ni nul,}} pas inv. sur $\mathcal{O}_X$\ldots

Formules

\[ \left\{\begin{aligned}\overline{x_1*x_2}&=\overline{x_1}*x_2=x_1*\overline{x_2}\\ \overline{x_1}*\overline{x_2}&=x_1*x_2\qquad\text{(résulte des précédentes)}\end{aligned}\right. \]

\deleted{i.e.} Prouvons $\overline{x_1}*x_2=\overline{x_1*x_2}$ il suffit pour $\unsure{x_2|x_1},x_2\in\Gamma P_{1}$

On a

\[ \begin{aligned}\overline{x_1}*x_2&=(T_1-x_1)*x_2=\underbrace{T_1}_{2x_1-b_1}*x_2-x_1*x_2=x_1*x_2-\\ &\hspace{18em}\deleted{\unsure{b}}b_1*x_2\\ \overline{x_1*x_2}&=T-x_1*x_2\deleted{\unsure{=}}=x_1*x_2-b_1\otimes b_2\end{aligned} \]

Or on a

\[ \begin{aligned}T&=2x_1*x_2-\underset{\substack{\Vert\\\deleted{2x_1-T_1}}}{b_1}\otimes\underset{\substack{\Vert\\\deleted{2x_2-\unsure{T_2|T_0}}}}{b_2}\qquad\text{d'où}\\ \overline{\unsure{a_1|x_1}*x_2}&=x_1*x_2-b_1\otimes b_2\qquad\text{OK.}\end{aligned} \]
% archive sheet number: 81
